% ECONOMICS_PROBLEM: {"id": "D82-513", "title": "Dimension dependence in learning a game by payments", "jel_code": "D82", "source_id": "OP-0155"}
% !TeX program = xelatex
\documentclass[11pt,a4paper]{article}
\usepackage[margin=23mm]{geometry}
\usepackage{fontspec}
\setmainfont{Latin Modern Roman}
\usepackage{amsmath,amssymb,mathtools}
\usepackage{xcolor,enumitem,booktabs,longtable,array,xurl,hyperref}
\definecolor{muted}{HTML}{53636C}
\hypersetup{colorlinks=true,urlcolor=blue,linkcolor=blue}
\setlength{\parindent}{0pt}
\setlength{\parskip}{5pt}
\newcommand{\R}{\mathbb R}
\newcommand{\E}{\mathbb E}
\newcommand{\N}{\mathbb N}
\newcommand{\ind}{\mathbf 1}
\newcommand{\Var}{\operatorname{Var}}
\newcommand{\Cov}{\operatorname{Cov}}
\newcommand{\supp}{\operatorname{supp}}
\DeclareMathOperator*{\argmin}{arg\,min}
\DeclareMathOperator*{\argmax}{arg\,max}
\newcommand{\entrymeta}[3]{{\sffamily\small\color{muted}\textbf{\texttt{#1}}\quad|\quad #2\par #3\par}}
\newcommand{\Problemref}[1]{\texttt{#1}}
\newcommand{\JELref}[2]{\texttt{#2} (JEL \texttt{#1})}
\begin{document}
\section*{Dimension dependence in learning a game by payments}
\textbf{Problem ID:} \texttt{D82-513}\quad \textbf{JEL:} \texttt{D82}\par
% BEGIN ECONOMICS STATEMENT
\label{problem:OP-0155}
% GAME_THEORY_PROBLEM: {"local_id": "GT-025", "original_number": 25, "kind": "R", "entry_kind": "statement_only", "published": false, "review_required": true, "category": "game_theory"}
\label{game:GT-025}
{\sffamily\small\color{muted}
{\sffamily\small\color{muted}\textbf{GT-025} | Learning in games and incentives\par R | Minimax specialization of an explicit research agenda\par}
\par}
\subsection*{Definitions and mathematical statement}
Let $N=\{1,\ldots,n\}$, $A=\prod_i A_i$, $m_i=|A_i|\geq2$, and $M=\prod_i m_i$. The unknown game has payoffs $U_i:A\to[0,1]$. In each of $T$ rounds a principal supplies each agent a context $s_i^t$ and chooses transfers $P_i^t:A_i\to[0,2]$. The agent chooses a mixed action $x_i^t$ in response to its context; the principal observes the joint mixed profile $x^t$. The agent's resulting utility is $U_i(x^t)+P_i^t(x_i^t)$; the agent is not given the transfer table as an additional decision signal.

The assumed learner guarantee is contextual external regret: with $\widehat U_i^t(a_i,x_{-i}^t)=U_i(a_i,x_{-i}^t)+P_i^t(a_i)$, for every context $s$ and prefix $t\leq T$,
\[
 \max_{a_i\in A_i}\sum_{\substack{q\leq t\, :\,s_i^q=s}}
 \left(\widehat U_i^q(a_i,x_{-i}^q)-\widehat U_i^q(x^q)\right)
 \leq C\sqrt T.
\]
Because opponent-dependent additive terms are strategically irrelevant, measure estimation error by
\[
 d(U,\widetilde U)=\inf_{(W_i)}\max_{i,a\in A}
 \left|U_i(a)+W_i(a_{-i})-\widetilde U_i(a)\right|,
 \qquad W_i:A_{-i}\to\mathbb R.
\]
For fixed action counts, $C$, and $\varepsilon>0$, let $L^*$ be the least horizon at which some allowed payment/context protocol guarantees $d(U,\widetilde U)\leq\varepsilon$ for every bounded game and every collection of agents satisfying the stated guarantee. Determine matching upper and lower dependence of $L^*$ on $(m_i)$, or on $M$ under an explicitly fixed family of action counts.
\subsection*{Short English statement}
How does the necessary number of rounds grow with game size when payments are used to learn strategically relevant utilities?
\subsection*{Variants and scope boundaries}
$M$ counts joint action profiles, not the sum of action counts. Observation of mixed strategies is stronger than observation of sampled actions. Round complexity, query complexity, and total payments are distinct objectives and require separate definitions. A probabilistic learner guarantee requires an additional confidence parameter.
\subsection*{Source relationship and status}
[S1], Section 6, explicitly leaves polynomial game-size factors unoptimized. The minimax quantity above is an editorial organization of that gap; no particular unmatched exponent is invented.
\subsection*{References}
\begingroup\footnotesize\raggedright
\textbf{[S1]} Brian Hu Zhang, Tao Lin, Yiling Chen and Tuomas Sandholm, \emph{Learning a Game by Paying the Agents}, arXiv:2503.01976v3 (2026 revision of a 2025 preprint). Locators: Sections 2--3, equations (2)--(3); Section 4, Theorems 4.2--4.4; Section 6, first future-work item. \url{https://arxiv.org/html/2503.01976v3}
\par\endgroup

\subsection*{Common conventions: Source mathematical conventions}

\textbf{Finite games.} For a finite set $A$, $\Delta(A)$ is its probability
simplex. Players choose independent mixed strategies in Nash equilibrium;
correlation is allowed only when explicitly part of the solution concept.
In a game with payoff functions $u_i$ and strategy profile $x$, an additive
$\varepsilon$ Nash equilibrium satisfies
\[
 u_i(x)\geq u_i(x_i',x_{-i})-\varepsilon
 \quad\text{for every player $i$ and every admissible deviation $x_i'$}.
\]
For numerical approximation, payoffs are normalized as locally specified.
An $\varepsilon$ payoff residual is distinct from distance at most
$\varepsilon$ to the set of exact equilibria. Point convergence, convergence
of empirical play, and convergence in distance to an equilibrium set are
also distinct.

\textbf{Computational representations.} Unless stated otherwise, a complexity
question takes rational data with binary encoding; polynomial time is measured
in total bit length. A fixed-accuracy polynomial algorithm is not necessarily
polynomial in $1/\varepsilon$, and neither is necessarily polynomial in
$\log(1/\varepsilon)$. Succinct games and value, demand, or sampling oracles
must declare their interfaces. Exact real solutions require an explicit
representation; an unspecified list of arbitrary real numbers is not a
finite computational output. A claimed algorithmic barrier is meaningful
only for its specified model and reductions.

\textbf{Dynamic games.} Histories, information, observation, and allowable
randomization are part of the model. A stationary strategy depends only on
the current state; a positional strategy in a deterministic graph game
chooses one action at each controlled vertex. Nash equilibrium at an initial
state and equilibrium after every history are different requirements.
An expectation of a pathwise $\liminf$ payoff, a $\liminf$ of expected averages,
a limiting expected average, and a uniform finite-horizon equilibrium are
different criteria. None is silently substituted for another.

\textbf{Allocation and mechanisms.} A complete allocation assigns every item
exactly once unless otherwise stated. Zero-valued goods, negative values,
capacity constraints, feasibility, lotteries, and copies of goods affect
fairness definitions and are specified locally. Dominant-strategy incentive
compatibility, Bayesian incentive compatibility, universal truthfulness,
and truthfulness in expectation impose different inequalities. Whenever
an objective ratio has zero denominator, its convention is stated or those
instances are excluded.

\textbf{Infinite-dimensional models.} Expectations, integrals, stochastic
equations, and optimization objectives require their local regularity,
integrability, filtrations, and admissible domains. A backward--forward
mean-field system is a boundary-value problem; it is not an initial-value
semiflow without an additional construction. Long-time, large-population,
and vanishing-noise limits require an order of limits and a convergence
topology. If a minimum need not be attained, the objective is its infimum
or supremum and the approximation requirement is stated separately.
% END ECONOMICS STATEMENT
\end{document}
